\documentclass[12pt,a4paper,oneside]{article}
\usepackage[brazil]{babel}
\usepackage[utf8]{inputenc}
\usepackage{olymp}

\setcounter{problem}{0}

\begin{document}

\begin{problem}{Diagonais}{diagonais}{2}

Em álgebra linear, as diagonais de uma matriz quadrada são os conjuntos de elementos da matriz que unem dois cantos opostos. Existem duas diagonais: a diagonal principal e a diagonal secundária. Veja o exemplo:

\begin{center}
\includegraphics[scale=0.5]{images/exercise_38_0.png}
\end{center}

Nesta matriz, os elementos da diagonal principal são os interceptados pela seta vermelha (1, 5 e 9) e os da diagonal secundária são os interceptados pela seta azul (3, 5 e 7).

Faça um programa para ler uma matriz quadrada de tamanho $N \times N$ e escrever os elementos de suas duas diagonais.

%\Notes
%
%Observações, se tiver.

\InputFile

A primeira linha contém um inteiro $N$, que indica o tamanho da matriz. Em seguida virão $N$ linhas com $N$ inteiros cada, indicando os elementos da matriz. Restrição: $1 \leq N \leq 1000$.

\OutputFile

Seu programa deve gerar $2N+1$ linhas de saída. As primeiras $N$ linhas devem conter os elementos da diagonal principal, um em cada linha. Em seguida, uma linha em branco, e novamente $N$ linhas, cada uma com um valor, desta vez contendo os elementos da diagonal secundária. Escreva os elementos sempre de cima para baixo.

% \Notes
% Preste atenção aos tipos das variáveis, pois $F(90) \approx 2^{61}$.

\Example

\begin{example}
\exmp{
3
1 2 3
4 5 6
7 8 9
}{
1
5
9\\

3
5
7
}%
\end{example}

\begin{example}
\exmp{
4
5 6 7 8
4 5 6 7
3 4 5 6
2 3 4 5
}{
5
5
5
5\\

8
6
4
2
}%
\end{example}

%Comentários sobre o exemplo, se necessário.
% Este exemplo corresponde ao da imagem mostrada acima.
\end{problem}

\end{document}
